\documentclass[10pt]{article}
\usepackage[margin=0.55in, top=0.45in, bottom=0.45in]{geometry}
\usepackage{enumitem}
\usepackage{titlesec}
\usepackage{hyperref}
\usepackage{tabularx}
\usepackage{booktabs}
\usepackage{xcolor}

\titlespacing*{\section}{0pt}{5pt}{2pt}
\titlespacing*{\subsection}{0pt}{3pt}{1pt}
\titleformat{\section}{\normalfont\bfseries\normalsize}{\thesection.}{0.4em}{}
\titleformat{\subsection}{\normalfont\bfseries\small}{\thesubsection.}{0.4em}{}
\setlength{\parskip}{2pt}
\setlength{\parindent}{0pt}
\pagestyle{empty}

\begin{document}

\begin{center}
{\Large \textbf{PhysID: Physical Property Fingerprinting via Manipulation}}\\[2pt]
{\normalsize 16-662 Robot Autonomy --- Project Proposal}
\end{center}

\section{Problem Statement}

Given a set of visually identical cubes with different hidden physical properties (mass, friction, stiffness, center of mass), can a Franka Emika Panda autonomously identify each cube by performing simple physical experiments and then sort them into labeled bins?

Robots rely heavily on vision, but many task-critical properties are invisible. We use \textit{manipulation as perception}.

\section{Approach}

\subsection{Test Objects}
We prepare 4--6 red cubes (5\,cm, identical appearance) with distinct physical profiles:

\begin{center}
\small
\begin{tabularx}{0.95\textwidth}{lX}
\toprule
\textbf{Cube} & \textbf{Hidden Property} \\
\midrule
A & Light (50\,g), rigid, low friction (Teflon-coated) \\
B & Heavy (300\,g), rigid, high friction (rubber-coated) \\
C & Medium (150\,g), hollow/off-center CoM, rigid \\
D & Medium (150\,g), deformable (foam core), uniform density \\
E & Heavy (300\,g), rigid, low friction (metal-filled, Teflon) \\
\bottomrule
\end{tabularx}
\end{center}

\subsection{Experiments the Robot Executes}

Each experiment is a scripted motion primitive executed via FrankaPy/MoveIt:

\begin{enumerate}[nosep,leftmargin=*]
    \item \textbf{Weigh} --- Grasp, lift to fixed height, hold still 2\,s. Compute mass from joint torque residuals: $m = \tau_{\text{res}} / (J^T \cdot g)$. Separates 50\,g / 150\,g / 300\,g classes.
    \item \textbf{Slide} --- Push cube across table at constant velocity (5\,cm/s). Measure steady-state force $\Rightarrow$ $\mu = F_{\text{push}} / (mg)$. Distinguishes Teflon ($\sim$0.04) vs rubber ($\sim$0.8).
    \item \textbf{Squeeze} --- Close gripper with ramping force (0--20\,N over 3\,s). Record force vs.\ displacement. Rigid cubes: near-zero displacement; foam cube: measurable compression $\Rightarrow$ stiffness $k = \Delta F / \Delta x$.
    \item \textbf{Shake} --- Sinusoidal wrist oscillation ($\pm$30\textdegree, 2\,Hz) per axis. Fit inertia model via least-squares on Newton-Euler equations. Off-center CoM produces asymmetric torque.
    \item \textbf{Tilt} --- Rotate wrist 90\textdegree\ slowly. Torque-vs-angle profile deviates from uniform model if CoM is off-center. Cross-validates with shake.
\end{enumerate}

\subsection{System Architecture}

\textbf{Perception} (AprilTags for cube pose) $\rightarrow$ \textbf{Experiment Planner} (Bayesian information-gain selector --- picks the most discriminating test next) $\rightarrow$ \textbf{Motion Executor} (FrankaPy) $\rightarrow$ \textbf{Estimator} (physics-based parameter fitting from F/T data) $\rightarrow$ \textbf{Classifier} (assigns cube ID when confidence $>$\,95\%) $\rightarrow$ \textbf{Sort into labeled bins}.

The planner updates a belief distribution after each experiment and selects the action that maximally reduces remaining uncertainty. At most 3 experiments per cube.

\section{Evaluation}
\begin{itemize}[nosep,leftmargin=*]
    \item \textbf{Classification accuracy}: confusion matrix over 20+ trials (target: $>$90\%)
    \item \textbf{Information-gain vs.\ random}: compare experiment count needed to reach 95\% confidence
    \item \textbf{Ablation}: which single experiment is most/least discriminating
\end{itemize}

\section{Key References}
\begin{enumerate}[nosep,leftmargin=*]
    \small
    \item Kruzliak et al., ``Interactive Learning of Physical Object Properties Through Robot Manipulation,'' IROS 2024. \href{https://arxiv.org/abs/2404.07344}{arXiv:2404.07344}
    \item Chen et al., ``Learning Object Properties Using Robot Proprioception via Differentiable Robot-Object Interaction,'' ICRA 2025. \href{https://arxiv.org/abs/2410.03920}{arXiv:2410.03920}
    \item Lambert, ``Identifying Objects' Inertial Parameters with Robotic Manipulation,'' MIT MEng 2023. \href{https://dspace.mit.edu/handle/1721.1/151518}{MIT DSpace}
\end{enumerate}

\end{document}
